\chapter{Literature Review}
\label{ch:literature}

The quantitative prediction of corporate bankruptcy has moved through three
overlapping phases, from classical statistical models over the adoption of
machine learning to a growing concern with class imbalance and model
interpretability. This chapter reviews the literature most relevant to 
the present work and isolates the gap this thesis addresses.

The study of bankruptcy prediction initially involved \textcite{beaver1966},
who examined financial ratios one at a time to separate failing from healthy
firms. \textcite{altman1968} then combined five ratios into a single
discriminant function, the Z-score, and later \textcite{ohlson1980} advanced
the field by introducing logit models. These methods are extensively discussed in the
literature reviewed in this thesis and are therefore only briefly noted here.
What matters for the present work is their common limitation. Although these
models are transparent and have coefficients that can be easily understood in
economic terms, they assume a fixed functional form and are thus incapable of
dealing with potential non-linear interactions in the data. The study by
\textcite{liang2016} is representative of the shift towards machine learning and
is also the origin of this dataset. At the same time, \textcite[562]{liang2016}
forced the two classes into an even balance by simple stratified sampling, a
decision that has since been questioned and is a design choice this 
thesis also does not follow.

Since bankruptcies are rare, the treatment of class imbalance has turned into a
major methodological issue. \textcite{veganzones2018} demonstrate that standard
models tend to perform badly as soon as bankrupt firms start to become too rare
in the training data, which means that some kind of rebalancing is required.
There are two types of solutions and they differ according to the kind of data
the model ends up learning. The oversampling methods, of which SMOTE
\parencite{chawla2002} is the most commonly used, create additional synthetic
observations for the minority class until the classes are balanced. In contrast,
specific undersampling methods reduce the number of cases in the majority class, so
that each observation on which the model is trained actually corresponds to a
firm that really existed. This difference is the methodological commitment of
the thesis. The closest study in this area is \textcite{wangliu2021}, in which
they examine several undersampling methods on exactly the dataset used here,
combined with repeated stratified cross-validation and the $F_2$-score as the
primary metric. The best setup that they achieve, using edited nearest
neighbours with a Naive Bayes classifier, attains an $F_2$ score of
$0.423$ \parencite[12]{wangliu2021}. The sampling approach, the idea of
leakage-aware validation and the choice of metric form the base that this thesis
extends. Cluster-based undersampling in particular has already existing work
outside of \textcite{wangliu2021}, since \textcite[11]{ansahnarh2024} also keep the
cluster centroids of the majority class on the same data, making use of the same
Python implementation as is used here, even though it is part of a more complex
overall setup.

There is also a large amount of more recent research which finds much higher
performance by combining synthetic oversampling with flexible learning methods,
an approach that \textcite{liashenko2023} and \textcite{amirshahi2024} survey
across both balanced and imbalanced settings. \textcite[305]{pham2025}, who also
provide a general overview of sampling strategies, achieve accuracies close to
$98.6\%$ using SMOTE and an artificial neural network, while \textcite{suresh2025}
and \textcite{sinap2025} report similar results with similar approaches.
Nevertheless, high accuracy under these setups has to be read with care. 
Some studies balance the data before splitting it, so the models are partly 
evaluated on synthetic firms, but even when the test set instead keeps the 
natural $96.8/3.2$ class distribution, accuracy alone is misleading because a 
model that predicts every firm as non-bankrupt already achieves $96.8\%$ accuracy.
This thesis takes these studies as points of differentiation rather than as 
templates and accepts the lower reported numbers in exchange for training only 
on real firms.

The final aspect of the issue relates to the features, the financial ratios,
which are responsible for the predictions. Some studies address this by carrying
out a selection at an early stage, either by selecting the features before
training, for instance \textcite{veganzones2018} who eliminate highly correlated
ratios and \textcite{chaudhary2023} who applies a Pearson-correlation screen, or
by transforming the features, as \textcite{ko2024} and \textcite{sinap2025} do
by means of factor analysis and principal components. Even though these
filtering methods are simple, they evaluate the features before and
independently of the model, whereas this thesis is concerned with finding
out what the trained model in fact relies on. In this regard, post-hoc
importance methods have been used, as \textcite{ansahnarh2024} report
random-forest importances and \textcite{sinap2025} is the only one of the
studies looked at here that makes use of SHAP, even though each of them provides
only a single importance measure without checking whether other measures would
be in agreement. Together, the framework being trained entirely on real firms
through undersampling and the comparison of importance methods make up the gap that 
this thesis aims to address. The technical
background of these methods is given in
Chapters~\ref{ch:methodology} and~\ref{ch:featureimportance}.


% highlight more the contribution of the thesis